\documentclass[11pt,letterpaper]{article}
\usepackage[T1]{fontenc}
\usepackage[utf8]{inputenc}
\usepackage{lmodern}
\usepackage[margin=1in]{geometry}
\usepackage{amsmath,amssymb,amsthm,mathtools}
\usepackage{booktabs,array,longtable}
\usepackage{enumitem}
\usepackage{microtype}
\usepackage{fancyhdr}
\usepackage{listings}
\usepackage{xcolor}
\usepackage{xurl}
\usepackage[hidelinks]{hyperref}
\usepackage[numbers,sort&compress]{natbib}
\usepackage{needspace}
\setlength{\headheight}{14pt}
\pagestyle{fancy}
\fancyhf{}
\fancyhead[L]{\small ProveIt / Unknot Recognition}
\fancyhead[R]{\small Certified sparse incidence}
\fancyfoot[C]{\thepage}
\setlength{\parskip}{0.35em}
\setlength{\parindent}{0pt}
\setlist{itemsep=0.25em,topsep=0.35em}
\lstset{basicstyle=\small\ttfamily,breaklines=true,columns=fullflexible,
 keepspaces=true,frame=single,rulecolor=\color{black!30},aboveskip=1em,belowskip=1em}
\newtheorem{theorem}{Theorem}[section]
\newtheorem{lemma}[theorem]{Lemma}
\newtheorem{proposition}[theorem]{Proposition}
\newtheorem{corollary}[theorem]{Corollary}
\newtheorem{definition}[theorem]{Definition}
\newtheorem{remark}[theorem]{Remark}
\newtheorem{example}[theorem]{Example}
\newcommand{\NN}{\mathbb N}
\newcommand{\ZZ}{\mathbb Z}
\newcommand{\supp}{\operatorname{supp}}
\newcommand{\poly}{\operatorname{poly}}
\newcommand{\bits}{\operatorname{bits}}
\newcommand{\calO}{\mathcal O}
\newcommand{\code}[1]{\nolinkurl{#1}}
\newcommand{\UU}{[r]}
\newcommand{\emptysetweight}{h(\varnothing)}
\newcommand{\ms}{\mathrm{ms}}
\title{\textbf{From Dense Port Tables to Polynomial Sparse Incidence}\\[0.5em]
\large Certified orbit-query acceleration for compressed unknot hierarchies}
\author{Research contribution prepared for Vladimir Reshetnikov's ProveIt project}
\date{8 October 2026}
\begin{document}
\maketitle
\begin{abstract}
The maintained marked-component observer in ProveIt's unknot-recognition
program uses exact interval-pairing orbit counts, canonical marked-union reuse,
and a dense Boolean M\"obius transform. Its output has $2^r$ entries for $r$
ports even when only a few component signatures occur. We replace that dense
representation by an exact positive sparse-zeta reconstruction, using the
existing unweighted Agol--Hass--Thurston (AHT) kernel without expanding a
single interval point.

For $s$ nonempty signatures, recovery requires $O(sr)$ oracle queries. Balanced
block deletion sharpens this to
$O(1+s+\sum_T |T|\log(1+r/|T|))$.
An independent certificate checks each positive signature through its immediate
proper subsets and a final mass identity, rather than re-running discovery or
enumerating all subsets. Crucially, for interval-union ports, polynomial support
size already follows from the classical weighted AHT theorem. Consequently the
sparse observer and its certificates have polynomial bit complexity in the
encoded input, with no logarithmic restriction on the number of ports. This is
a certified unweighted-query realization of a consequence of existing weighted
orbit theory, not a new general weighted-orbit complexity theorem.

We also recover incidence classified by a supplied binary orientation character;
the lifted histogram has exactly the same support, so its reconstruction needs
no second support search. A standard-library implementation, two byte-identical
pinned upstream kernels, independent literal-graph audits, adversarial tests,
and paired measurements accompany the proofs. Several sparse families improve
substantially; dense and small controls exhibit overhead. The result removes a
specific exponential-format bottleneck. It neither constructs a complete
unknot hierarchy nor proves a general quasi-polynomial unknot recognizer.
\end{abstract}

\vspace{0.5em}
\textbf{Status.} Complete written proofs of the reduction, extraction,
certificate, and binary-cover statements; a polynomial-complexity corollary
explicitly dependent on classical AHT. Executable exact-arithmetic validation;
no Lean formalization or first-in-literature priority claim. The delivered
code is an additive research package, not a committed production change.
\newpage
\tableofcontents
\newpage

\section{Research objective and the inspected bottleneck}
\label{sec:scope}

\subsection{Why this is the next component to improve}
The requested objective is to continue improving the exact recognizer in
\code{Topology/UnknotRecognition}, aiming ultimately at quasi-polynomial
asymptotic complexity. The inspected sources distinguish a complete recognition
pipeline from local filters, compressed group operations, normal-surface
observers, and hierarchy primitives. That distinction is essential: improving
one exact primitive is useful, but does not establish a bound for a composition
whose other stages or number of calls remain uncontrolled.

The most recent maintained orbit work is a particularly suitable target.
It supplies an unweighted AHT count kernel, local proof replay, and a marked
incidence observer. The synthesis explains that repeated unions are cached,
but that the histogram and its proof-index table still have $2^r$ slots
\cite{upstream,orbitreport}. Thus the representation itself can be exponentially
larger than its nonzero information. This paper changes the representation,
the discovery algorithm, and the certificate contract together.

The concrete outcome is a function taking a supplied compressed interval
relation and marked interval unions and returning the \emph{complete sparse
histogram} of the component signatures. The result includes multiplicity,
including the multiplicity of the unmarked signature. It is not merely a
sample of components or a list of signatures encountered before a resource
limit.

\subsection{Snapshot and reproducibility boundary}
The inspected repository ref is
\begin{center}
\small\code{7518823550fbc8c217bc8be0113fe52002e77465}.
\end{center}
Two modules are delivered unchanged, with their Git blob identities checked
before they are loaded:
\begin{center}
\small
\begin{tabular}{@{}ll@{}}
\toprule
Module & Git blob\\
\midrule
\code{interval_orbits.py} & \code{e86050aadbcef6f17fdddff9cbd1abc0e7ca18e8}\\
\code{interval_orbit_verify.py} & \code{0ccb56a0e8b5f1255384121d7417441314727720}\\
\bottomrule
\end{tabular}
\end{center}
The second module imports neither the orbit producer nor its scheduling
helpers. Its proof checker is reused as a trusted local-relation layer; the
new sparse certificate verifier does not call the sparse producer either.

All delivered measurements use the pinned count module with its
\code{periodic_rule='aht'} setting, retaining the classical sufficient merger
criterion. The sharper Fine--Wilf implementation remains unchanged in the
vendored file, but its acceleration is not needed for this paper's complexity
argument or benchmark comparison. This choice avoids making the new argument
depend on a separate analysis of a modified orbit scheduler.

The dense comparison arm is a small research ablation: identical coning,
canonical-union caching, the same pinned classical kernel, and dense M\"obius
inversion. It is \emph{not} a byte-for-byte copy of the full maintained incidence
module or recognition dispatch. No full upstream test suite, Regina integration,
native normal-surface corpus, or whole-knot benchmark is claimed to have been
run for this contribution.

\subsection{The strongest result, and its attribution}
The important point is stronger than a conditional statement for small $s$.
AHT's weighted algorithm returns compressed blocks of equal orbit-weight
vectors in polynomial time \cite[Section 6, Theorem 16]{aht}. Marking each port
with one nonnegative coordinate shows that the number of distinct signatures
is already polynomial in the interval-encoded input. The new unweighted-query
reconstruction inherits that polynomial bound.

Accordingly, the contribution is not that weighted orbit information has
suddenly become polynomial-time computable. It is an explicit, independently
certifiable route from the already maintained \emph{unweighted} engine to a
sparse incidence interface, avoiding a second weighted-transfer implementation
and avoiding the maintained dense output format. The output-sensitive oracle
bound, the sparse local certificate, the support-reusing binary-cover extension,
and the measured implementation are the repository-facing developments.

\section{Exact input and output models}
\label{sec:model}

\begin{definition}[Interval relation]
Let $X_N=\{0,\ldots,N-1\}$, where $N\geq0$ is binary encoded. An interval pairing
is an isometry between two nonempty integer intervals of equal length. In the
supplied representation it is a row $(a,b,c,d,\epsilon)$ with inclusive
endpoints, $b-a=d-c$, and $\epsilon\in\{1,-1\}$. Its action is
\[
 x\longmapsto
 \begin{cases}
 x+c-a,&\epsilon=1,\\
 a+d-x,&\epsilon=-1.
 \end{cases}
\]
The relation $\sim$ is the equivalence relation generated by $k$ such pairings.
Its equivalence classes are called orbits or components.
\end{definition}

Pairings and their inverses generate the same equivalence relation. The input
normalizer places the domain interval no farther right than the range,
retaining the order of the supplied pairing list. Identity pairings and
repeated pairings are allowed. The empty universe requires an empty pairing
list. No input contains one record per point unless the caller chooses such
an expanded representation.

\begin{definition}[Ports and signatures]
A port $A_i\subseteq X_N$, $1\leq i\leq r$, is supplied as a finite union of
half-open intervals $[u,v)$. Endpoints satisfy $0\leq u\leq v\leq N$.
Within each port, overlap and adjacency are merged and empty intervals are
removed. Ports may overlap, coincide, or be empty.
For a component $O$, put
\[
 \sigma(O)=\{i\in[r]:O\cap A_i\neq\varnothing\},\qquad
 h(T)=\#\{O:\sigma(O)=T\}.
\]
The total component count is $C=\sum_{T\subseteq[r]}h(T)$.
\end{definition}

Let $M$ be the number of marked intervals in the input, or any upper bound
on the number after normalization. Let
\[
 B=\max\{1,\lceil\log_2(N+1)\rceil\},\qquad
 s=\#\{T\neq\varnothing:h(T)>0\},\qquad
 s_0=\#\{T:h(T)>0\}.
\]
Thus $s_0$ is either $s$ or $s+1$. The output is a sorted list of the $s_0$
positive pairs $(T,h(T))$, with a support encoded as an $r$-bit mask. An absent
entry means exactly zero. The empty output means $C=0$; it does not mean an
unfinished computation.

Counts satisfy $0\leq h(T)\leq C\leq N$, so every count uses $O(B)$ bits.
The output needs $O(s_0(r+B))$ bits, apart from ordinary serialization
conventions. A dense array instead has $2^r$ slots even for one positive
signature. These are different output contracts, not interchangeable
representations with an identical mandatory cost.

\subsection{Three sizes that must not be conflated}
The number of points $N$, the number of components $C$, and the number of
signature types $s_0$ can differ by enormous factors. For a static universe
with no pairings, $C=N$. If every point has the same nonempty signature, then
$s_0=1$. Both $N$ and $C$ can equal $2^{16000}$ while the multiplicity is
represented by one binary integer. A valid sparse algorithm must remove a
whole signature's multiplicity in one step; an algorithm that removes one
component at a time would lose the intended compressed complexity.

\section{From marked unions to positive zeta queries}
\label{sec:oracle}

\subsection{Coning without expanding a marked interval}
For a selected union $A\subseteq X_N$, define
\[
 H(A)=\#\{O:O\cap A\neq\varnothing\}.
\]
Suppose the canonical union is nonempty and consists of $m_A$ disjoint,
nonadjacent half-open intervals. Choose its first point as an anchor. For an
interval $[u,v)$ of length at least two, add the pairing
\[
 [u,v-2]\longrightarrow[u+1,v-1].
\]
This is one inclusive interval pairing, not $v-u-1$ separately stored edges.
For every interval other than the first, add one singleton pairing between
its first point and the anchor. At most $2m_A-1\leq 2M$ pairings are added.
For a length-one first interval there is no internal translation.

\begin{lemma}[Coning identity]
\label{lem:cone}
Let $C_A$ be the component count after the preceding relations are added to
the original interval system. If $A\neq\varnothing$, then
\[
 H(A)=C-C_A+1.
\]
For the empty marked union, $H(\varnothing)=0$ without an additional query.
\end{lemma}
\begin{proof}
The internal translation identifies consecutive points in each marked
interval. The singleton anchor relations then identify all marked points.
Every original component meeting $A$ therefore joins the same new component.
No added relation touches an original component disjoint from $A$, so those
components remain distinct and unchanged. Exactly $H(A)$ old components
become one, giving $C_A=C-H(A)+1$.
\end{proof}

The argument applies to overlapping ports, disconnected original systems,
static points, and orientation-reversing original pairings. Coning is used
only to ask a question about the original relation; the resulting merged
system is not substituted for that relation in later geometry.

\subsection{The Boolean zeta function}
For $S,U\subseteq[r]$, put
\[
 F(S)=H\!\left(\bigcup_{i\in S}A_i\right),\qquad
 G(U)=C-F([r]\setminus U).
\]

\begin{lemma}[Positive downset representation]
\label{lem:zeta}
\[
 G(U)=\sum_{T\subseteq U}h(T).
\]
In particular, $G([r])=C$, $G(\varnothing)=h(\varnothing)$, and $G$ is
monotone under inclusion.
\end{lemma}
\begin{proof}
A component is excluded from $F([r]\setminus U)$ precisely when it meets none
of the ports outside $U$. This is equivalent to its signature being a subset
of $U$. Counting those components proves the formula. Monotonicity follows
because every summand is nonnegative.
\end{proof}

For a nonempty actual union of the excluded ports, Lemma~\ref{lem:cone} gives
the useful direct formula
\begin{equation}
 \label{eq:direct}
 G(U)=C_A-1.
\end{equation}
When that actual union is empty, $G(U)=C$. A nonempty mask can describe an
empty union if all its selected ports are empty; testing the mask alone is
therefore insufficient.

\subsection{Caching keys that preserve the exact question}
The oracle stores the baseline count once. It also stores completed $G(U)$
values by mask and coned counts by the \emph{canonical actual union} of the
excluded ports. Distinct masks with the same union produce exactly the same
coning rows in the same order, so a count and its proof can safely be reused.
Equality of two counts is never used as evidence that two unions are
interchangeable.

This retains the maintained observer's useful sharing mechanism. What changes
is that only the masks needed by discovery or certification are constructed.
A cache indexed by all $2^r$ masks in advance would reintroduce the bottleneck.

\section{Positive sparse reconstruction}
\label{sec:recovery}

\subsection{Residuals and minimal positive supports}
Let $\mathcal K$ be a set of signatures whose multiplicities have already
been recovered exactly. Define
\begin{equation}
\label{eq:residual}
 R_{\mathcal K}(U)=G(U)-\sum_{T\in\mathcal K,\,T\subseteq U}h(T).
\end{equation}
Under the stated invariant,
\[
 R_{\mathcal K}(U)=\sum_{T\notin\mathcal K,\,T\subseteq U}h(T).
\]
Thus the residual is again a nonnegative, monotone zeta function. In
particular, a zero residual rules out every unrecovered signature in the
queried downset; it is not merely a cancellation among positive and negative
terms.

\begin{lemma}[Minimal positive extraction]
\label{lem:minimal}
If $S$ is inclusion-minimal among sets with $R_{\mathcal K}(S)>0$, then
$S\notin\mathcal K$, $h(S)>0$, and $R_{\mathcal K}(S)=h(S)$.
\end{lemma}
\begin{proof}
Every unrecovered proper subset $T\subsetneq S$ lies inside some
$S\setminus\{a\}$, $a\in S\setminus T$. Minimality says that the residual
there is zero, so nonnegativity forces $h(T)=0$. The only remaining possible
positive summand at $S$ is $h(S)$. Since the residual is positive, that summand
exists and has not already been removed.
\end{proof}

The multiplicity is therefore recovered exactly in one subtraction. It may
be exponential in $B$ as an ordinary integer magnitude, but its bit length is
only $O(B)$.

\subsection{Flat deletion algorithm}
First obtain $C=G([r])$ and $h(\varnothing)=G(\varnothing)$. Retain the empty
signature if its multiplicity is positive. While the sum of recovered
multiplicities is smaller than $C$, start with $S=[r]$ and inspect the ports
in a fixed order. Delete $a$ from $S$ whenever
$R_{\mathcal K}(S\setminus\{a\})>0$. At the end, recover
$h(S)=R_{\mathcal K}(S)$ and repeat.

\begin{theorem}[Exact flat recovery]
\label{thm:flat}
For every nonnegative Boolean histogram, the preceding algorithm returns the
complete sparse histogram after exactly $s$ nonempty-support extraction
iterations. It needs at most $2+s(r+1)$ logical zeta requests, before caching,
and does not require advance knowledge of $s$ or of the support sizes.
\end{theorem}
\begin{proof}
At the start of each iteration, the residual at $[r]$ is positive because
some total mass remains. Each successful deletion preserves that property.
Suppose a deletion of $a$ failed at an intermediate set $S'$. The final set
$S$ is contained in $S'$, so
$S\setminus\{a\}\subseteq S'\setminus\{a\}$. Monotonicity keeps its
residual zero. Hence every retained element fails a final single-element
deletion, and every proper subset of $S$ is contained in one such deletion.
The final $S$ is minimal positive. Lemma~\ref{lem:minimal} therefore recovers
one exact new nonempty signature.

The empty signature has already been removed, so this iteration cannot
return it. An iteration subtracts the full positive mass of one previously
unknown signature and never changes any other multiplicity. Exactly $s$
iterations suffice. When the recovered mass is $C$, no missing positive
entry can remain. The initial two requests, at most $r$ deletion tests, and
one final weight request per iteration give the stated bound.
\end{proof}

The implementation computes the remaining total by an integer mass
accumulator; it does not re-query the full set at each iteration. A mask cache
can also eliminate repeated requests for the empty set or a final support.

\subsection{Balanced block deletion}
\label{sec:balanced}
Deleting one port at a time wastes queries when an unknown signature is
small. Partition the ordered $r$ ports into a balanced binary interval tree.
At a current set $S$, visit a tree block $D$ and test
$R_{\mathcal K}(S\setminus D)$. If positive, delete the entire block and do
not visit its children. If zero and the block is not a leaf, recursively
visit its two children in their fixed order. A failed leaf is retained.
Start by visiting the root block; its empty-set query is normally already
cached.

\begin{theorem}[Support-sensitive query bound]
\label{thm:balanced}
Balanced block deletion returns a minimal positive support. If that support
has size $d\geq1$, its number of block tests is at most
\begin{equation}
 \label{eq:blockbound}
 b(r,d)=1+2\sum_{\ell=0}^{\lceil\log_2 r\rceil-1}\min(2^\ell,d)
 =O\!\left(1+d\log\left(1+\frac rd\right)\right).
\end{equation}
Consequently the entire histogram is recovered using
\begin{equation}
 \label{eq:sparsequeries}
 Q=O\!\left(1+s+\sum_{T\neq\varnothing:\,h(T)>0}
 |T|\log\left(1+\frac r{|T|}\right)\right)
\end{equation}
zeta requests, and in particular $Q=O(1+sr)$.
For $r=0$, only the baseline/empty count is needed.
\end{theorem}
\begin{proof}
Every successful block deletion preserves positive residual. Every retained
leaf has failed its deletion test, and the monotonicity argument from
Theorem~\ref{thm:flat} keeps that failure valid as subsequent deletions shrink
$S$. Thus the final set is minimal positive.

Fix this final support $T$. Whenever a block $D$ fails at an intermediate
set $S'$, all residual signatures contained in $S'$ meet $D$. The final $T$
is one such signature, so $D\cap T\neq\varnothing$. At depth $\ell$, the
tree blocks are disjoint. At most $\min(2^\ell,d)$ blocks at that depth can
fail. Every visited block other than the root is a child of a failed
nonleaf. There are therefore at most twice that many visits at the next
depth. Summing gives~\eqref{eq:blockbound}; an unbalanced last level only
reduces the count.

To estimate the sum, levels with $2^\ell\leq d$ contribute $O(d)$ in total,
and the remaining $O(1+\log(r/d))$ levels contribute $O(d)$ each. One additional
request obtains the weight of the final support. Add the initial requests
and sum over all extraction iterations to obtain~\eqref{eq:sparsequeries}.
\end{proof}

\begin{remark}[Order influences cost, not validity]
The final minimal signature and the number of requests can depend on the
port order. Every fixed order is correct, and the displayed bound does not
assume a favorable ordering. Finding an order minimizing expensive coned
queries is a separate optimization problem. A query count alone does not
measure the differing cost of its underlying interval systems.
\end{remark}

\subsection{A useful information lower bound}
\begin{proposition}[Single-signature lower bound]
\label{prop:lower}
Suppose an oracle has one unknown positive signature $T$ of known size $d$
and known weight one, on $r$ ports. Any deterministic adaptive reconstruction
algorithm using zeta queries needs at least
$\lceil\log_2\binom rd\rceil$ queries in its worst case.
\end{proposition}
\begin{proof}
Every query returns either zero or one, according as $T\not\subseteq U$ or
$T\subseteq U$. A depth-$q$ binary decision tree has at most $2^q$ leaves,
whereas the $\binom rd$ possible supports must be distinguished.
\end{proof}

For sparse ranks away from the degenerate dense endpoint, this gives the
same familiar $d\log(r/d)$ information scale as balanced deletion. It is not
an optimality claim for every $d$, for multiple weighted signatures, or for
special interval geometry. For example, if $d=r$ is promised in advance,
the support is already known; our general-purpose algorithm does not use
that promise.

Hidden weighted-hypergraph reconstruction is an established query-learning
subject. Bshouty and Mazzawi obtain substantially more refined nonadaptive
bounds for constant-rank weighted hypergraphs \cite{bshouty}. Here the
requirements are different: arbitrary support sizes, simple exact adaptive
queries, compatibility with compressed interval geometry, and an independent
certificate for the complete returned histogram. No general query-optimality
or priority claim is made.

\section{Sparse certificates without dense forward sums}
\label{sec:certificates}

\subsection{The algebraic certificate criterion}
Let proposed positive entries be $(T_1,w_1),\ldots,(T_t,w_t)$, with distinct
nonempty masks and positive integer weights. A canonical certificate orders
them first by cardinality, then by mask. Let $w_0\geq0$ be the proposed empty
multiplicity, and define
\[
 K_j(U)=w_0+\sum_{i<j:\,T_i\subseteq U}w_i.
\]
The $w_0$ term occurs in every downset, including the empty downset.

\begin{theorem}[Local sparse certificate]
\label{thm:certificate}
Assume $G$ is the zeta transform of a nonnegative histogram of total mass $C$.
The proposed entries are exactly that histogram if the following equalities
hold:
\begin{align}
 G(\varnothing)&=w_0,\label{eq:certempty}\\
 G(T_j)&=K_j(T_j)+w_j &&(1\leq j\leq t),\label{eq:certweight}\\
 G(T_j\setminus\{a\})&=K_j(T_j\setminus\{a\})
 &&(a\in T_j,\ 1\leq j\leq t),\label{eq:certdelete}\\
 C&=w_0+\sum_{j=1}^t w_j.\label{eq:certmass}
\end{align}
Including the baseline, at most
$2+t+\sum_j|T_j|$ orbit-count queries are sufficient before sharing of
coincident unions and empty-union shortcuts.
\end{theorem}
\begin{proof}
Equation~\eqref{eq:certempty} proves the empty multiplicity. Inductively assume
the earlier nonempty entries have also been proved correct. Subtracting those
entries leaves an actual nonnegative residual histogram.
Equation~\eqref{eq:certdelete} says that its zeta value vanishes on every
immediate proper subset of $T_j$. Every proper subset is contained in one
of these immediate subsets. Nonnegativity therefore forces every residual
proper-subset weight to be zero. Equation~\eqref{eq:certweight} then proves
that the actual residual weight at $T_j$ is precisely $w_j$.

After all entries have been checked, every proposed multiplicity has been
proved exact. The remaining mass is nonnegative, and
\eqref{eq:certmass} says that it is zero. Hence there are no omitted positive
signatures anywhere, not just inside the tested downsets. The query bound
counts one baseline, one empty-set value, and the support and immediate
proper-subset values for every nonempty entry.
\end{proof}

This certificate does not require the verifier to know which supports the
producer tested, which block order it chose, or why it stopped searching.
It verifies the answer, rather than the execution strategy. The final mass
identity is indispensable: local correctness of several entries alone says
nothing about signatures outside their downsets.

\subsection{Why deletion witnesses cannot simply be discarded}
Consider a proposed histogram with a single weight-one support $[r]$. Its
total mass and full-set value agree with the histogram having support
$[r]\setminus\{a\}$ for any fixed $a$, provided $r\geq2$. Among all subsets
$U\subseteq[r]$, the corresponding zeta functions differ only at
$U=[r]\setminus\{a\}$. Thus a general exact verifier for this claimed
full-support histogram must distinguish each such alternative. It needs all
$r$ immediate-deletion values in this oracle model. A final mass check plus
one value per output row is not, in general, a sound replacement.

\subsection{Binding algebraic values to supplied geometry}
The delivered certificate has the schema \code{sparse-incidence-v1}. It
contains the canonical input universe, ordered pairing rows, ordered normalized
ports, a baseline local orbit proof, the sparse entries, and a list of coned
union proofs. The producer first recovers the histogram, then obtains any
additional values required by Theorem~\ref{thm:certificate}. Those additional
calls use the \emph{same remaining} query and cycle allowances. It emits only
proofs needed by the certificate, not proofs of discovery-only queries.

The new verifier performs the following checks. It validates and compares
the complete input binding, rejecting noncanonical certificate geometry and
Boolean values masquerading as integer fields. It replays the baseline proof
with the independent upstream local checker. For each requested zeta value,
it reconstructs the actual excluded-port union with an independent endpoint
sweep, rebuilds the coning rows, and replays the bound orbit proof. It then
checks~\eqref{eq:certempty}--\eqref{eq:certmass}. Missing, duplicate, unused,
or misbound union proofs are rejected. Repeated equal unions are checked
once, but never equated merely because their claimed counts match.

The verifier imports no reconstruction algorithm or orbit producer. It shares
only public geometry normalization with the producer. In the tests, the orbit
producer is deliberately disabled after proof construction and the completed
sparse proof still verifies.

\subsection{What the local orbit proofs establish}
The reused upstream verifier recognizes identity deletion, static contraction,
reflection trimming, periodic merger, transmission by a certified power of a
partial inverse, and exposed suffix truncation \cite{orbitreport}. It checks
those relations locally instead of insisting on the producer's preferred
scheduling choices. A proof must finish with an empty universe and account
for the claimed number of removed singleton orbits. An unfinished valid prefix
is not an orbit-count certificate.

The algebraic certificate theorem depends only on truthful query values, and
the coning lemma binds those values to the original components. Thus the proof
layers have separate obligations: local relation replay proves counts,
coning proves marked-union interpretation, and nonnegative sparse algebra proves
completeness of the histogram. None of these layers establishes that the
supplied interval system is the correct knot-exterior model of a diagram.

\section{Bit complexity and polynomial support}
\label{sec:complexity}

\subsection{Output-sensitive accounting}
Let $P(K,B)$ bound a completed classical unweighted AHT call, including optional
recording when needed, on $K$ pairings and $B$-bit endpoints. Let $V(K,B)$ and
$L(K,B)$ bound independent replay time and serialized binary trace size for
such produced proofs. These bounds are polynomial by classical AHT and the
local trace construction; precise optimized exponents are not needed here.
Logs in these bounds may be replaced by $\max(1,\log)$ to cover zero and
unit-sized inputs.

Every coned query has $K\leq k+2M$ pairings on the \emph{same} universe.
No coning operation increases endpoint bit length. A straightforward sorted
union costs $O(M\log(M+1)B)$ bit operations, with additional mask operations
polynomial in $r$. The delivered residual calculation scans the list of
previous positive entries. For each pair it performs an $r$-bit containment
test and, when needed, a $B$-bit subtraction.

\begin{theorem}[A conservative sparse bit bound]
\label{thm:bits}
With $Q$ as in~\eqref{eq:sparsequeries}, histogram discovery can be implemented
in bit time
\begin{equation}
\label{eq:bitbound}
 O\!\left(
 Q\left[P(k+2M,B)+M\log(M+1)B+(s+1)(r+B)\right]
 \right)
 +\poly(k+M+r+B).
\end{equation}
The analogous certificate construction adds at most
$Q_c=2+s+\sum_T|T|$ potential count queries and polynomial sparse bookkeeping.
Verification replaces $P$ by $V$ for its at most $Q_c$ distinct proofs.
No $2^r$ term is required in discovery, certificate output, or replay.
\end{theorem}
\begin{proof}
Use Lemma~\ref{lem:cone} for every required marked-union value and charge each
actual count by $P(k+2M,B)$. Building the corresponding union and coning rows
has the stated sorting and linear-output cost. For a residual value, at most
$s+1$ already recovered masks and counts are inspected. All residual values
are between zero and $C$, so their integers have $O(B)$ bits. Add the input
normalization, sparse output sorting, and initial allocations to the polynomial
term. The certificate query bound is Theorem~\ref{thm:certificate}; its
verifier uses only those masks and the supplied traces.
\end{proof}

The displayed estimate deliberately does not assume constant-time arithmetic
on arbitrarily large Python integers. Dictionary lookups in the implementation
use ordinary expected hash-table costs. A deterministic ordered dictionary
would add polynomial comparison factors, without changing the polynomial
conclusion. A still simpler list-based cache also suffices for that conclusion.

Proof storage is bounded by
\[
 O\bigl(Q_cL(k+2M,B)+Q_cMB+s_0(r+B)\bigr)
 +\poly(k+M+r+B).
\]
During a traced run the oracle can temporarily retain proofs of discovery
queries that will not be emitted; this uses the analogous bound with $Q+Q_c$.
The implementation does not promise a minimal-memory streaming certificate
producer. That is an optimization opportunity, not an unaccounted exponential
allocation.

\subsection{Why the sparse support is polynomial for interval inputs}
It would be incomplete to leave the geometry-independent parameter $s$
entirely unexplained. For arbitrary positive Boolean functions it may be
$2^r-1$. For the present interval-encoded inputs, classical weighted orbit
theory gives a stronger conclusion.

\begin{lemma}[Polynomial signature support from weighted AHT]
\label{lem:polysupport}
For $k$ interval pairings and $r$ interval-union ports with a total of $M$
intervals, the number $s_0$ of positive component signatures is bounded by a
polynomial in $k+M+r+B$.
\end{lemma}
\begin{proof}
The case $N=0$ is immediate. For $N>0$ and $r>0$, assign to each point the nonnegative vector
\[
 w(x)=\bigl(\mathbf1_{A_1}(x),\ldots,\mathbf1_{A_r}(x)\bigr)\in\ZZ_{\geq0}^r.
\]
The vector is constant on a common partition with at most $2M+1$ nonempty
cells, since changes can occur only at port endpoints. The partition and its
vectors have a polynomial-size representation, for example by an endpoint
sweep keeping an $r$-bit active-port mask. For a component $O$, its orbit-weight
vector is
\[
 v(O)=\sum_{x\in O}w(x).
\]
Its positive coordinates are exactly $\sigma(O)$.

Apply the weighted AHT theorem to this piecewise constant vector weight.
Its compressed output is a list of interval blocks, each block representing
many components sharing one orbit-weight vector \cite[Section 6]{aht}.
The number of written blocks is polynomial in the input parameters. The weight
dimension is $r$, the number of input cells is at most $2M+1$, and a safe
numerical weight bound is $D=1+rN$, so $\log D=O(B+\log(r+1))$. A signature
can occur only as the positive-coordinate support of one of the vectors on
those output blocks. The number of distinct such supports is at most the
number of blocks, and hence is polynomial. When $r=0$, there is at most the
empty signature and the conclusion is immediate.
\end{proof}

The proof uses the \emph{compressed output format}, not an impossible explicit
list of all $C$ components. A block may stand for $2^{16000}$ components. One
must also use a coordinate per supplied port, not one coordinate per point.
The theorem's polynomial dependence on weight dimension is precisely what
permits polynomially many ports.

\begin{theorem}[Polynomial sparse interval-incidence observer]
\label{thm:polynomial}
Given binary-encoded interval pairings and interval-union ports, the complete
sparse signature histogram can be computed and furnished with independently
replayable certificates in polynomial bit time in the encoded input size.
There is no requirement that $r=O(\log n)$, $r=O((\log n)^2)$, or that component
multiplicities be polynomially bounded as ordinary numbers.
\end{theorem}
\begin{proof}
Lemma~\ref{lem:polysupport} gives polynomial $s_0$. The number of sparse
requests in Theorem~\ref{thm:balanced} is at most $O(1+sr)$; the certificate
request bound is of the same polynomial order. Classical unweighted AHT gives
polynomial $P,V,L$ for the coned instances. Substitute these bounds into
Theorem~\ref{thm:bits} and the certificate-storage estimate.
\end{proof}

This theorem is a consequence of established weighted AHT combined with the
explicit unweighted-query reduction. It should not be advertised as a new
general complexity theorem for weighted orbit extraction. Its practical
advantage is that the executable path uses the already available unweighted
proof engine. A direct weighted implementation may have better constants and
would additionally return coordinate magnitudes, not just their supports.

\subsection{Dense output remains an exponential requirement}
If a caller insists on an array with an addressable slot for every subset,
writing that array costs $\Omega(2^r)$ slots regardless of the preceding theorem.
Even a zero universe with $r$ empty ports has a dense histogram of that length.
Thus the theorem changes the output interface rather than claiming to fill an
exponentially long object in polynomial time.

The distinction also explains the dense benchmark control. Giving every one
of the $2^r$ signatures a static point is possible, but the corresponding
bit-pattern interval marks have exponentially many interval pieces. Such a
family is not a counterexample to polynomiality in $k+M+r+B$.

\section{Incidence refined by a binary orientation character}
\label{sec:signed}

\subsection{Order reversal and orientation parity are different inputs}
Attach to each generating pairing an additional bit $p\in\{0,1\}$. This bit
records the supplied character, not necessarily the sign of the interval
isometry. The double cover has points $(x,e)$, $e\in\{0,1\}$, with relations
\[
 (x,e)\sim(g(x),e+p).
\]
It is encoded in two consecutive blocks of $N$ points, using $2k$ interval
pairings. Each port is lifted to both sheets, so the number of marked interval
pieces is at most $2M$. This encoding doubles the size numerically but adds
only one endpoint bit.

A base component is \emph{consistent} if every closed walk in its relation
graph has total parity zero. If the supplied bits are the actual surface
orientation character, consistency means orientability. Without that provenance,
only the binary-label consistency statement is justified.

\begin{lemma}[Lift support is unchanged]
\label{lem:lift}
Every connected component of the double cover projects surjectively onto one
base component and has its exact port signature. A consistent base component
has two lifted components; an inconsistent one has one. Consequently, if $b(T)$
and $\ell(T)$ are the base and lifted histograms, respectively, then
\[
 \supp b=\supp\ell,\qquad
 b(T)\leq\ell(T)\leq2b(T).
\]
\end{lemma}
\begin{proof}
Lift a path from a chosen base point to any other point in its component.
Starting at a chosen sheet determines a lifted endpoint, proving surjectivity
of each lifted component. Because a port is lifted on both sheets, a lifted
component meets that port exactly when its base component does.

If all closed walks have parity zero, assigning a sheet label by path parity
is path-independent. The two possible initial labels give two disjoint
components. If a closed walk has parity one, its lift joins the two sheets
above one point, and path lifting then joins them above every point. The cover
is connected over that base component. Summing by signature proves the bounds.
\end{proof}

\begin{theorem}[Typed sparse incidence and support reuse]
\label{thm:typed}
Let $o(T)$ and $u(T)$ count consistent and inconsistent base components of
signature $T$. Then
\[
 o(T)=\ell(T)-b(T),\qquad u(T)=2b(T)-\ell(T).
\]
After the base sparse support has been recovered, the lifted histogram can be
recovered with one baseline and at most $s_0$ additional zeta requests, before
caching. No second balanced support search is necessary. Typed incidence has
polynomial bit complexity for interval-union marks.
\end{theorem}
\begin{proof}
Lemma~\ref{lem:lift} gives $b=o+u$ and $\ell=2o+u$, whose solution is the stated
formula. Order the known support by increasing cardinality. For a support $T$
in that order, every positive proper-subset support is already known and
processed. Since the lifted support equals the base support, its weight is
\[
 \ell(T)=G_{\mathrm{lift}}(T)-
   \sum_{S\subsetneq T:\,b(S)>0}\ell(S).
\]
This triangular calculation uses one value per known support. The total lifted
count supplies a final consistency check. For independent certification, the
ordinary sparse certificate criterion is applied to both histograms; it may
require additional deletion values, all under the shared budget. Polynomiality
follows from Theorem~\ref{thm:polynomial} on the doubled encoding.
\end{proof}

The implementation's signed result contains rows $(T,o(T),u(T))$, preserving
zeros in the last two fields when the other field is positive. Its proof
contains separately replayable base and cover certificates. The independent
signed verifier reconstructs the cover from the supplied signed pairings,
checks both certificates, checks equality of supports, and recomputes the two
nonnegative differences. A genuine character-changing input mutation cannot
be hidden by retaining the old cover proof.

\section{Counterexamples to overstrong conclusions}
\label{sec:barriers}

\subsection{Signed weights invalidate zero-based extraction}
For the abstract signed histogram
\[
 h(\{1\})=1,\qquad h(\{2\})=-1,
\]
the total zeta value at $\{1,2\}$ is zero although the histogram is not zero.
A total-mass stopping rule would fail. This is why the reconstruction theorem
is about \emph{nonnegative component counts}; it must not be transferred
unchanged to Euler-characteristic coefficients, signed polynomial coefficients,
or arbitrary integer-weight cancellation data. In Section~\ref{sec:signed},
subtraction is performed only after two genuine nonnegative histograms have
been recovered and verified.

\subsection{Equal incidence is not equal future attachment behavior}
\begin{proposition}[Four-point attachment obstruction]
\label{prop:attachment}
There exist two interval relations with identical complete port-incidence
histograms that acquire different component counts after the same new
pointwise attachments.
\end{proposition}
\begin{proof}
Take four points $a_1,a_2,b_1,b_2$ and ports
$A=\{a_1,a_2\}$ and $B=\{b_1,b_2\}$. In the first relation $E$, identify
$a_1\sim b_1$ and $a_2\sim b_2$. In the second relation $E'$, identify
$a_1\sim b_2$ and $a_2\sim b_1$. Both histograms consist of the single row
$h(\{A,B\})=2$. Add to each relation the same two attachments
$a_1\sim b_1$ and $a_2\sim b_2$. The first relation still has two components;
the second becomes connected. These are valid interval inputs by assigning
the four consecutive coordinates $0,1,2,3$ and using singleton pairings.
\end{proof}

Thus a signature histogram is an exact observer, not a complete state for
arbitrary future gluing. It loses the association between individually named
points, as well as multiplicity at a port, cyclic order, slope, and boundary
attachment maps. A hierarchy algorithm may use the histogram to answer the
specific incidence question, but cannot replace the original boundary state
by that histogram without a separate congruence theorem for its allowed
attachments.

\subsection{Polynomial local work does not bound hierarchy search}
A complete recognition bound also needs an algorithm for producing the
supplied interval relations from the input diagram, controlling the encoding
of cut surfaces and marked boundaries, generating the necessary candidates,
and stopping after a bounded number of stages. None of those facts follows
from an incidence query theorem.

Lackenby's 2026 hierarchy preprint supplies important geometric algorithms,
but its Section 9 explicitly distinguishes counts of steps from fully
specified running-time estimates and leaves parameters to be bounded there
\cite{lackenby}. Treating that discussion as an already available complete
quasi-polynomial implementation would skip precisely the cost obligations
that the repository is trying to resolve.

\section{A precise interface to a future quasi-polynomial hierarchy}
\label{sec:hierarchy}

\begin{proposition}[Composition accounting]
\label{prop:composition}
Suppose a proposed hierarchy procedure on an $n$-crossing input makes at most
$H(n)$ sparse-incidence calls. Suppose every supplied incidence instance has
encoded length at most $L(n)$, and that all other work, including candidate
production, port extraction, certificate replay, and terminal tests, takes
at most $W(n)$ bit operations. Then its total work is bounded by
\[
 W(n)+H(n)L(n)^{O(1)}.
\]
If $W,H,L$ are each at most $2^{(\log(n+2))^{O(1)}}$, the resulting bound is
quasi-polynomial. This conclusion assumes correctness and completeness of the
underlying hierarchy procedure; it does not supply them.
\end{proposition}
\begin{proof}
Theorem~\ref{thm:polynomial} gives a fixed polynomial bound for each incidence
call, including its chosen proof operations. Sum over the $H(n)$ calls and
add the remaining work. Products and fixed powers of the displayed
quasi-polynomial scales remain in that scale.
\end{proof}

The point is not the elementary summation. It is that the incidence observer
no longer contributes an independent $2^{r(n)}$ factor. Polynomially many
ports are permissible when their interval encodings are polynomially large.
For the project's more specific target $n^{O(\log n)}=2^{O((\log n)^2)}$,
one needs compatible $O((\log n)^2)$ exponents for the \emph{whole} expression,
not merely membership in the broader quasi-polynomial class.

Depth and number of calls are also different quantities. A search tree of
depth $O(\log n)$ and branching factor $n^{O(1)}$ has quasi-polynomially many
nodes, not necessarily polynomially many. If branching or encoded surface
sizes grow faster, a logarithmic depth slogan does not repair the bound.
The interval-incidence theorem removes one obstacle; it does not replace
this global accounting.

\section{Implementation contracts and resource behavior}
\label{sec:implementation}

\subsection{Module separation}
The source package has a small mathematical core. \code{reconstruct.py}
implements flat and balanced positive-zeta extraction independently of
interval geometry. \code{geometry.py} validates and normalizes public input.
\code{oracle.py} performs canonical union queries through the unchanged AHT
kernel. \code{verify.py} checks sparse certificates without importing discovery.
\code{signed.py} implements the binary cover, support reuse, and typed replay.
\code{api.py} provides the ordinary sparse interface and the dense ablation.

The kernel loader checks the exact Git blob hash of each vendored source file
before importing it. The producer and verifier are loaded separately. This
protects the claimed source identity and makes an accidental edited fixture
visible; it is not a claim that source hashing proves the algorithm correct.

\subsection{Completion, interruption, and mutable state}
A successful query returns \code{status='COMPLETE'}, its full sparse histogram,
and an optional certificate. A query interrupted by its support, total-query,
or total-AHT-cycle allowance returns \code{INCONCLUSIVE} with no histogram,
no total-count assertion, and no certificate. It may retain work statistics.
There is no conversion from resource exhaustion to ``unknot'' or ``knotted.''

The total cycle budget is shared among the baseline, all coned calls, and
additional certification calls. The signed wrapper shares it between the base
and the doubled system as well. The same is true for the total query allowance.
AHT cycles are structural events, not a wall-clock or bit-operation meter.
The optional callback provides cooperative cancellation, and its exceptions
propagate rather than being mislabeled as a successful mathematical result.

All core objects are constructed from normalized copies. The supplied input
lists are not mutated. A complete sparse answer produced before a failed
certificate stage is not exposed as a fully certified answer by the combined
API; the combined operation returns inconclusive instead. A caller needing
separate uncertified and certified products can use separate explicit calls
with separate documented budgets, rather than silently resetting one budget.

\subsection{Serialization and hostile-input scope}
Endpoint integers and component counts may exceed Python's default limit for
decimal-string conversion. The transport layer uses exact hexadecimal strings
for sufficiently large integers, without modifying the interpreter's global
limit. Boolean values are not accepted as mathematical integers. The proof
format carries no floating-point quantity used in a correctness decision.

Polynomial verification in the explicit proof length does not mean that any
untrusted JSON file is cheap to parse or store. Callers should impose limits
on serialized bytes, nesting, and wall time before accepting external
certificates. The delivered verifier is a mathematical checker, not a sandbox
for arbitrarily large hostile inputs.

\section{Exact validation and controlled measurements}
\label{sec:experiments}

\subsection{Completed validation}
The following checks were completed in this environment. The standard-library
suite contains 33 test methods, including invalid input, resource-boundary,
proof-mutation, and noninterference tests. Its exhaustive subtest checks all
$3^8=6561$ histograms with eight possible supports on three ports and weights
in $\{0,1,2\}$, using both flat and balanced extraction. Three hundred additional
positive histograms exercise larger masks and multiplicities. Tests also check
the explicit balanced-tree bound and the logarithmic-query behavior of a
single hidden singleton on 1024 ports.

A separate seeded audit generates 1000 interval systems with up to 50 points
and up to ten ports. It computes components independently by literal union-find
and compares every sparse entry. It also assigns random binary parities and
checks the typed result against an independent graph traversal with parity
labels. Every supplied base and cover certificate is replayed. The audit retains
all inputs, expected sparse outputs, and statistics in
\code{results/validation.json}. All 1000 ordinary comparisons and all 1000
binary-character comparisons agree, with no failed proof replays.

The ordinary certificates in that audit exercise every supported local AHT
rule: deletion, contraction, trimming, merger, transmission, and truncation.
The exact event counts are in the audit file. This is finite evidence about
the implementation, not a replacement for the general proofs. The test log
also includes a $2^{16000}$-point transport/replay example and a large parallel
chain family without point expansion.

Adversarial tests alter multiplicities, omit supports, duplicate entries,
change port bindings, truncate local traces, delete or duplicate union proofs,
insert an unused otherwise valid proof, and replace an integer with a Boolean.
They check a specific case in which certificate construction needs one more
orbit query than discovery: the shared allowance stops it rather than
resetting. The four-point attachment obstruction is an executable regression.

\subsection{Benchmark design}
Nine supplied interval workloads are measured with five arms: cached dense
inversion, an identical dense A/A control, flat sparse recovery, balanced sparse
recovery, and balanced recovery with certificate production and independent
replay included in its timing. Each arm starts with fresh oracle state. There
is one warmup and seven measured rounds per workload, and arm order is shuffled
with seed 261008491. Thus the principal experiment contains 45 warmup and 315
measured calls. Result comparison is outside the timed region; replay is
inside the timed region only for the explicitly labeled proof/replay arm.

The implementation uses CPython; its full version and platform string are
recorded in \code{results/benchmark.json}, along with source hashes, every
sample, arm order, query count, and proof size. Timings are observations from
this container, not machine-independent complexity claims. The table reports
median milliseconds and the median of the paired dense/balanced ratios. A
paired ratio need not equal the ratio of the two separately displayed medians.

\begin{table}[htbp]
\centering\small
\begin{tabular}{@{}lrrrrr@{}}
\toprule
Workload & Dense & Flat & Balanced & Proof/replay & Ratio\\
\midrule
Static, coincident & 0.683 & 0.054 & 0.073 & 0.180 & 9.25\\
Static, nested & 0.763 & 0.300 & 0.369 & 0.884 & 2.05\\
Static, disjoint & 20.130 & 1.579 & 0.783 & 1.638 & 25.44\\
8-block chain, coincident & 0.560 & 0.450 & 0.476 & 0.682 & 1.21\\
8-block chain, nested & 1.879 & 1.837 & 1.799 & 2.456 & 1.01\\
8-block chain, disjoint & 22.849 & 6.541 & 4.039 & 5.389 & 5.70\\
16-block chain, disjoint & 380.408 & 47.020 & 25.777 & 28.915 & 14.70\\
All 32 signatures & 3.134 & 3.477 & 3.877 & 6.824 & 0.86\\
Small mixed relation & 0.294 & 0.244 & 0.267 & 0.589 & 1.10\\
\bottomrule
\end{tabular}

\caption{Completed same-kernel component measurements. Times are milliseconds.
``Ratio'' is paired dense time divided by balanced sparse time; values below
one indicate a slowdown. No row is a whole-knot recognition benchmark.}
\label{tab:timings}
\end{table}

The static workloads have universe size $2^{4000}$ and eight ports. The chain
workloads link consecutive equal blocks by translations. Their base orbits are
the common offsets within a block, so a block length of $2^{500}$ represents
that many parallel components. The eight-block chains use six ports; the
sixteen-block chain uses eight ports and block length $2^{128}$. Marks are
explicit coordinate intervals with coincident, nested, or disjoint layouts.
They are not asserted to have been extracted from a native knot-exterior
attachment construction.

\begin{table}[htbp]
\centering\small
\begin{tabular}{@{}lrrrrr@{}}
\toprule
Workload & $r$ & $s_0$ & Dense calls & Sparse calls & Proof KiB\\
\midrule
Static, coincident & 8 & 2 & 2 & 2 & 26.4\\
Static, nested & 8 & 8 & 9 & 9 & 89.2\\
Static, disjoint & 8 & 8 & 256 & 21 & 179.6\\
8-block chain, coincident & 6 & 2 & 2 & 2 & 18.4\\
8-block chain, nested & 6 & 6 & 7 & 7 & 52.1\\
8-block chain, disjoint & 6 & 6 & 64 & 14 & 66.2\\
16-block chain, disjoint & 8 & 8 & 256 & 21 & 48.8\\
All 32 signatures & 5 & 32 & 32 & 32 & 19.0\\
Small mixed relation & 3 & 4 & 8 & 8 & 3.3\\
\bottomrule
\end{tabular}

\caption{Actual orbit calls, including the baseline. Here $s_0$ includes a
positive empty signature. Proof size is the JSON byte count divided by 1024;
it is a transport measurement, not a claim of minimal binary encoding.}
\label{tab:queries}
\end{table}

\subsection{Interpretation, including regressions}
The disjoint sparse cases are the clearest demonstration of the mechanism.
Dense inversion needs every distinct subset union, whereas balanced discovery
asks a small selection of them. The eight-port disjoint examples reduce 256
count calls to 21. The six-port disjoint chain reduces 64 to 14. These are exact
call counts, independent of timing noise.

Coincident and nested marks illustrate a different effect. Canonical-union
caching already reduces the number of underlying count calls, so sparse
recovery cannot improve those counts much or at all. It instead avoids dense
mask enumeration and inversion. On the eight-block nested chain, the result
is near a tie; certificate production and replay add visible cost. The
all-signatures control has 32 positive rows and 32 count calls in either
representation. Sparse bookkeeping adds overhead there. The small mixed
case is also too small and noisy to justify a broad speed claim.

The A/A samples are retained specifically to expose this scale of variation.
No aggregate claim that the complete recognizer has accelerated is inferred
from these tables. The observer is not called by ordinary knot dispatch in the
inspected pipeline, and the dense baseline is a controlled ablation rather than
a rerun of that pipeline.

An exploratory classical-kernel benchmark containing a 32-block chain was
stopped by this session's execution-time limit before a complete result file
was produced. It contributes no numbers to the tables. The delivered benchmark
uses the explicitly named 16-block case and completed all its rounds. This
change concerns experimental tractability, not the theorem or the compressed
input model.

\subsection{Large-port scaling without a dense allocation}
\begin{table}[htbp]
\centering\small
\begin{tabular}{@{}lrrrrr@{}}
\toprule
Family & $r$ & Dense slots & Sparse rows & Calls & Time (ms)\\
\midrule
Disjoint static & 16 & $2^{16}$ & 16 & 49 & 1.610\\
Disjoint static & 32 & $2^{32}$ & 32 & 113 & 4.867\\
Disjoint static & 64 & $2^{64}$ & 64 & 257 & 15.431\\
Disjoint static & 128 & $2^{128}$ & 128 & 577 & 53.372\\
One active port & 64 & $2^{64}$ & 1 & 2 & 0.118\\
One active port & 256 & $2^{256}$ & 1 & 2 & 0.309\\
One active port & 1024 & $2^{1024}$ & 1 & 2 & 1.404\\
\bottomrule
\end{tabular}

\caption{Completed balanced-only scaling runs, seven repetitions each. The
``dense slots'' column is the output-format requirement, not a measured dense
run or an extrapolated timing.}
\label{tab:scaling}
\end{table}

For disjoint static ports, the universe has size $2^{500}$ and the ports
partition it into equal blocks for the listed powers of two. The exact
histogram has one singleton signature per port. For the one-active-port
family, the universe has size $2^{16000}$, exactly one port is the entire
universe, and the remaining ports are empty. The unique nonempty signature
has multiplicity $2^{16000}$. In the 1024-port case, balanced discovery uses
21 block tests and only two distinct underlying orbit counts after union
reuse. No dense array with $2^{1024}$ entries was constructed or timed.

\section{Integration plan and preservation of existing guarantees}
\label{sec:integration}

The package is designed as an additive research directory, for example
\code{Topology/UnknotRecognition/research/sparse_incidence_20261008/}.
It does not overwrite any maintained source and does not allocate a numbered
report slot that might already have been used in a concurrent repository
update. \code{INTEGRATION.md} gives the precise production-facing steps.

A safe first integration is an explicit sparse observer API, separate from
ordinary knot recognition. It should accept the same interval-pairing model,
normalize ports once, and return a tagged sparse histogram. Existing callers
expecting dense mask indexing must either remain on the dense interface or
request a deliberate, separately capped expansion. Silent shape changes are
not compatible with exact callers.

The certificate schema should remain distinct from the existing dense
incidence schema. A dense verifier's $2^r$ proof-index contract is not an
appropriate compatibility layer for the new sparse result. The proof layer
can reuse the maintained local orbit verifier after its exact version and
encoding contract are checked. The included source-probe script reports blob
matches and fails on mismatches rather than claiming to have validated an
unseen future revision.

Before promotion, run the full maintained tests and the native normal-arc
adapter audits on the intended checkout, including geometric provenance tests.
Measure sparse and dense observers on the same native inputs, retain A/A
controls, and keep whole-recognition timings separate. Until such a production
caller is introduced, the correct performance statement is a faster supplied
incidence query, not a faster complete recognizer.

\section{Further research questions and concrete next targets}
\label{sec:future}

\subsection{A weighted engine versus repeated unweighted queries}
The polynomial support theorem comes from weighted AHT, while the executable
method makes repeated unweighted calls. A directly certified vector-weight
transfer engine could compute all port counts, component Euler data, and
incidence in one pass. The concrete target is an independently replayable
weighted-transfer contract with controlled integer growth and a paired
comparison against the present reduction. The relevant question is not
whether a polynomial algorithm exists---that is classical---but whether one
can obtain a smaller practical proof and lower total cost within the
maintained architecture.

\subsection{Sharper support bounds tied to geometry}
The weighted theorem supplies a polynomial upper bound without a small
explicit exponent in this paper. Can normal-arc systems or actual hierarchy
ports force a much smaller bound, such as a near-linear number of signature
types in the number of pairing and marking intervals? A useful approach is to
count how many different supports can be created by weighted transfer and
static contraction, then test extremal examples generated by genuine surface
encodings. No such near-linear claim is assumed here.

\subsection{Residual-query data structures}
The delivered code scans every previously recovered support when evaluating
a residual. This gives a conservative $Q(s+1)(r+B)$ term. A dynamic sparse
subset-sum structure, compressed trie, or rank-separated index could reduce
that term on structured supports. The target must remain exact and should
be benchmarked against the list scan, which is inexpensive for very small
$s$. An index whose preprocessing enumerates the Boolean cube would defeat
the purpose.

\subsection{Cost-aware balanced trees and port equivalences}
Balanced deletion minimizes neither total coned interval count nor total AHT
cycles. Can one choose a binary partition tree using port overlap, repeated
unions, or measured query cost, while retaining a provable approximation
bound for an explicit cost model? Identical ports may be quotiented before
search, and their bits restored afterward. Any more aggressive equivalence
must prove equality of component-incidence behavior rather than merely
similarity of coordinate ranges.

\subsection{Certificate size and streaming replay}
Discovery and independent verification ask overlapping but not identical
sets of questions. Can a streaming producer retain just the proofs that will
be needed by the final sparse certificate without a costly restart? Can a
certificate verifier arrange the sparse equations to reduce repeated subset
arithmetic, while preserving the immediate-deletion information lower bound?
The implementation already rejects unused proof records, but it does not
minimize temporary traced storage.

\subsection{Finite covers beyond a binary character}
For a fixed finite group of order $q$, a regular cover similarly preserves
port signatures when every mark is lifted completely. A component lifts to
between one and $q$ components. The support-known triangular reconstruction
still applies. Determining detailed monodromy subgroup types would require
additional data or multiple covers; the two-by-two linear formulas of
Section~\ref{sec:signed} do not automatically generalize to identify them.
A precise next target is a small fixed-group family with a proved separating
set of covers and independently replayable subgroup-type counts.

\subsection{Attachment-compatible refinements}
Proposition~\ref{prop:attachment} shows that incidence alone is not a congruence
for pointwise attachments. What additional compressed data make it a congruence
for the \emph{particular} attachments occurring in a normal-surface hierarchy?
Candidate data include boundary multiplicities, ordered interval maps,
permutation actions on repeated sheets, and a finite set of slope labels.
The goal is a closure theorem: two summaries declared equivalent must remain
equivalent under every allowed continuation. This is more important than
showing that an isolated histogram is small.

\subsection{Certified port extraction and compressed cutting}
The local polynomial theorem requires ports given by polynomially many
binary intervals. A future geometric adapter must prove that this form is
preserved by cutting, crushing, and provenance updates. A successful adapter
should supply a certificate connecting every new marked interval to the
original triangulation and to the chosen normal coordinates, including
orientation characters. A compact normal-coordinate vector alone does not
prove that an arbitrary proposed collection of attachment ports is correct.

\subsection{Global potential and search accounting}
The decisive missing result is a potential controlling the number of hierarchy
states and the size of their complete encodings. A meaningful theorem would
bound $H(n)$ and $L(n)$ in Proposition~\ref{prop:composition}, and account for
candidate generation and terminal checks in $W(n)$. Geometric depth,
branching, compressed word length, and numerical surface weight should have
separate counters. The current contribution makes one of these costs
polynomial; it does not allow the remaining counters to be omitted.

\subsection{Native validation and formal proof boundaries}
The next experimental step is a native normal-surface corpus with adversarial
port arrangements, compared against both the maintained dense observer and
an independent geometric component extractor at feasible sizes. Whole-knot
measurements should begin only when a production caller actually uses the
new query. A tractable formalization target is the positive residual invariant
and Theorem~\ref{thm:certificate} over a finite Boolean lattice, separated from
AHT and geometric provenance. That would make the sparse completeness argument
machine-checkable without pretending to formalize an entire hierarchy at once.

\section{Conclusion}
The maintained dense marked-incidence observer admits an exact sparse
replacement that reuses certified unweighted interval-orbit counting.
Positive residuals locate minimal unknown signatures, whole multiplicities
are recovered at once, balanced deletion makes discovery sensitive to support
size, and local deletion equations plus total mass independently certify the
complete result. A binary cover preserves support and permits a second
histogram to be recovered without repeating support discovery.

The asymptotic conclusion for supplied interval data is polynomial, not merely
conditional on a logarithmic port bound. Its geometric support argument is an
explicit corollary of classical weighted AHT, duly attributed. The observed
speedups concern exactly this observer and are accompanied by dense and small
controls showing overhead. For general unknot recognition, the remaining
research burden is a certified attachment-compatible hierarchy with controlled
candidate generation, encoded sizes, and total work. The delivered primitive
removes a concrete exponential-format cost from that program without
concealing those unresolved obligations.

\appendix
\section{Executable usage and artifact inventory}
\label{sec:artifacts}
From the extracted package root:
\begin{lstlisting}
export PYTHONPATH="$PWD/src${PYTHONPATH:+:$PYTHONPATH}"
python -m unittest discover -s tests -v
python experiments/validate.py
python experiments/benchmark.py --rounds 7
python experiments/make_tables.py
python -m sparse_incidence analyze examples/mixed.json \
    --output results/mixed.json --certificate
python -m sparse_incidence verify examples/mixed.json results/mixed.json
cd paper && sh build.sh
\end{lstlisting}
No network connection is required. Python 3.10 or later is the intended minimum;
the actual tested version is recorded in the results. Rebuilding the article
requires a LaTeX installation providing the packages in its preamble and BibTeX.
The example commands analyze an interval relation, not a knot diagram.

\begin{longtable}{@{}p{0.40\linewidth}p{0.55\linewidth}@{}}
\toprule
Path & Contents\\
\midrule
\endhead
\code{paper/article.tex}, \code{article.pdf} & Complete article and compiled PDF.\\
\code{paper/references.bib} & Primary literature and immutable repository references.\\
\code{src/sparse_incidence/} & Sparse extraction, geometry, cached oracle, independent replay, signed cover, CLI.\\
\code{vendor/} & Two unchanged upstream modules, separately pinned by Git blob and SHA-256.\\
\code{tests/test_sparse.py} & 33 unit-test methods with exhaustive, seeded, and adversarial subcases.\\
\code{experiments/} & Literal graph controls, seeded audit, fresh-state paired benchmark, table generation.\\
\code{results/} & Completed test log, all audit inputs/results, raw timing samples, summary logs.\\
\code{examples/} & Ordinary and huge-multiplicity inputs and independently replayed certificates.\\
\code{integration/} & Local source-identity probe; no remote writes or automatic integration.\\
\code{CLAIMS.md}, \code{INTEGRATION.md} & Scope ledger and production-review obligations.\\
\code{PROVENANCE.json} & Repository ref, inspected sources, source hashes, runtime and scope.\\
\code{SHA256SUMS} & Delivered-file integrity manifest, excluding itself.\\
\bottomrule
\end{longtable}

\section{Worked sparse certificate on a mixed relation}
The small mixed example has $N=10$, the translation $[0,4]\to[5,9]$, and
ports
\[
 A_1=[0,2),\quad A_2=[6,9),\quad A_3=[0,1)\cup[9,10).
\]
The five components are $\{i,i+5\}$, $0\leq i<5$. Their signatures give
\[
 h(\{2\})=2,\quad h(\{3\})=1,\quad
 h(\{1,2\})=1,\quad h(\{1,3\})=1,
\]
and all other entries vanish. In particular $h(\varnothing)=0$.
A cardinality-ordered certificate first proves the two singleton entries.
It then checks the weight at $\{1,2\}$ after subtracting the known singleton
mass $2$, and checks both deletion values. The entry at $\{1,3\}$ is
handled similarly after subtracting the singleton mass $1$. The final total
$2+1+1+1=5$ rules out every missing signature, including $\{2,3\}$ and
$\{1,2,3\}$, without a dense forward transform.

The standalone proof file contains full local AHT witnesses for its actual
coned unions. A reader can mutate the claimed row weights or bindings and
replay the file using the CLI. The sparse verifier does not trust the
\code{COMPLETE} label in a surrounding result; it verifies the certificate
against the separately supplied input.

\section{Claim ledger}
\begin{longtable}{@{}p{0.42\linewidth}p{0.53\linewidth}@{}}
\toprule
Claim & Status\\
\midrule
\endhead
Coned-union identity & Proved; existing maintained construction, reused and independently reconstructed in replay.\\
Positive sparse extraction & Proved for every nonnegative Boolean histogram; implemented in two strategies.\\
Support-sensitive balanced bound & Proved with an explicit binary-tree count; checked on exhaustive and seeded examples.\\
Sparse local completeness certificate & Proved; independently implemented and adversarially tested.\\
Polynomial interval-incidence complexity & Proved as a consequence of classical weighted and unweighted AHT plus the explicit sparse reduction. Not a new general weighted-orbit theorem.\\
Binary-character refinement & Proved and implemented, including known-support cover recovery and shared budgets.\\
Polynomial general unknot recognition & Not claimed.\\
General quasi-polynomial unknot recognition & Not established by this contribution.\\
Whole-knot practical acceleration & Not measured; no production recognition dispatch was changed.\\
Native surface-adapter validation & Not run here; explicitly required before production promotion.\\
Full upstream suite & Not run here. The 33-method delivered suite and 1000-case independent audits were run.\\
Formal machine-checked proofs & Not provided; written proofs and executable exact checks only.\\
Priority beyond the inspected repository & Not claimed; related weighted-orbit and hidden-hypergraph literature is acknowledged.\\
\bottomrule
\end{longtable}

\bibliographystyle{plainnat}
\bibliography{references}
\end{document}
